\section{Anhang}
% Pro Code-Datei ein Unterabschnitt: ein, zwei Sätze, wofür die Datei da ist,
% dann \codefile[angezeigter Name]{Pfad}{Hinweis}. Zuerst das Notebook,
% danach die Module.
% Notebook exportieren:  jupyter nbconvert --to script auswertung.ipynb
%   (--no-prompt entfernt die "# In[n]:"-Zeilen,
%    --TemplateExporter.exclude_markdown=True die Markdown-Zellen)
% Beispielinhalt: Unterabschnitte und Dateien in code/ können gelöscht werden.

\subsection{Auswertungsnotebook}\label{sec:notebook}
Das Notebook \texttt{auswertung.ipynb} führt die gesamte Auswertung durch und
erzeugt \cref{fig:d}. Markdown-Zellen erscheinen als Kommentare,
\texttt{\# In[n]:} markiert die Zellen.
\codefile[auswertung.ipynb]{code/auswertung.py}{Jupyter-Notebook, exportiert}

\subsection{Modul \texttt{fitting.py}}\label{sec:fitting}
Gewichtete Anpassung einer Ursprungsgeraden; wird in Zelle~3 des
Notebooks (\cref{sec:notebook}) verwendet.
\codefile{code/fitting.py}{Python-Modul}

\subsection{Versuchsprotokoll}\label{sec:protokoll}
Das handschriftliche Versuchsprotokoll ist auf den folgenden Seiten
eingescannt.
% Alle Seiten des Scans (A4), je eine pro Seite in voller Textbreite;
% pagecommand={} behält Kopf- und Fußzeile.
\includepdf[pages=-, width=\textwidth, pagecommand={}]{assets/placeholder-protocol}
